\subsection{NILPOTENT LIE GROUPS}

PROPOSITION 12. Let G be a finite-dimensional Lie group. For L(G) to be nilpotent, it is necessary and sufficient that G have a nilpotent open subgroup.

Suppose that \(\mathbf{G}\) has a nilpotent open subgroup \(\mathbf{G}_0\) . By the Corollaries to Propositions 4 and 6, no. 2, \(\mathcal{C}^i\mathrm{L}(G_0) = \{0\}\) for sufficiently large \(i\) . Hence \(\mathrm{L}(G_0) = \mathrm{L}(G)\) is nilpotent.

Suppose that \(\mathrm{L}\langle \mathrm{G}\rangle\) is nilpotent. If \(\mathbf{K} = \mathbf{R}\) or \(\mathbf{C}\) , the identity component \(\mathbf{G}_0\) of \(\mathbf{G}\) is nilpotent by the Corollary to Proposition 4, no. 2, and \(\mathbf{G}_0\) is open in \(\mathbf{G}\) . If \(\mathbf{K}\) is ultrametric, the Corollary to Proposition 6, no. 2, proves that there exist an open subgroup \(\mathbf{G}_1\) of \(\mathbf{G}\) , an integer \(i > 0\) and a neighbourhood \(\mathbf{V}\) of \(e\) in \(\mathbf{G}\) such that \(\mathrm{C}^i\mathrm{G}_1\cap \mathbf{V} = \{\varepsilon \}\) . Then, if \(\mathbf{G}_0\) is a sufficiently small subgroup of \(\mathbf{G}_1\) , \(\mathrm{C}^i\mathrm{G}_0\subset \mathrm{V}\) , hence \(\mathrm{C}^i\mathrm{G}_0 = \{\varepsilon \}\) and \(\mathbf{G}_0\) is nilpotent.

Let \(\mathfrak{g}\) be a nilpotent Lie algebra. The Hausdorff series H(X, Y) corresponding to \(\mathfrak{g}\) has only a finite number of non-zero terms and we know (Chapter II, § 6, no. 5, Remark 3) that the law of composition \((x, y) \mapsto H(x, y)\) defines a group structure on \(\mathfrak{g}\) . Suppose further that \(\mathfrak{g}\) is complete and normable. Clearly the law H is a continuous-polynomial (Differentiable and Analytic Manifolds, R, Appendix). Hence \(\mathfrak{g}\) , together with the law H, is a Lie group G, said to be associated with \(\mathfrak{g}\) . By §4, no. 2, Lemmas 2 and 3, \(\mathrm{L}(\mathrm{G}) = \mathfrak{g}\) . The identity mapping \(\phi\) of \(\mathfrak{g}\) into \(\mathrm{G}\) is an exponential mapping of \(\mathrm{G}\) such that \(\phi(\lambda x)\phi(\lambda' x) = \phi((\lambda + \lambda') x)\) for all \(x \in \mathfrak{g}\) , \(\lambda \in \mathrm{K}\) , \(\lambda' \in \mathrm{K}\) . Every Lie subalgebra \(\mathfrak{h}\) of \(\mathfrak{g}\) admitting a topological supplement is a Lie subgroup \(\mathrm{H}\) of \(\mathrm{G}\) and \(\mathrm{L}(\mathrm{H}) = \mathfrak{h}\) .

PROPOSITION 13. Let G be a finite-dimensional simply connected nilpotent Lie group over R or C.

\begin{enumerate}
\def\labelenumi{(\roman{enumi})}
\item
  \(\exp_{\mathbf{G}}\) is an isomorphism of the Lie group associated with \(\mathrm{L}(\mathrm{G})\) onto \(\mathbf{G}\) .
\item
  Every integral subgroup of \(\mathbf{G}\) is a simply connected Lie subgroup of \(\mathbf{G}\) .
\end{enumerate}

Let \(\mathfrak{g} = \mathrm{L}(\mathrm{G})\) , which is nilpotent (Proposition 12). As two simply connected Lie groups over \(\mathbf{R}\) or \(\mathbf{C}\) which have the same Lie algebra are isomorphic (§ 6, no. 3, Theorem 3 (ii)), it suffices to prove the proposition when \(\mathrm{G}\) is the group associated with \(\mathfrak{g}\) . Then (i) and (ii) follow from what was said before the proposition.

PROPOSITION 14. Let G be a finite-dimensional connected Lie group over R or C.

\begin{enumerate}
\def\labelenumi{(\roman{enumi})}
\item
  If G is nilpotent, \(\exp_{G}\) is étale and surjective.
\item
  If K = C and \(\exp_{G}\) is étale, then G is nilpotent.
\end{enumerate}

Let \(G'\) be the universal covering space of G. Let \(\phi\) be the canonical morphism of \(G'\) onto G. Then \(\exp_G = \phi \circ \exp_{G'}\) (§ 6, no. 4, Proposition 10) and hence (i) follows from Proposition 13 (i).

If \(\mathbf{K} = \mathbf{C}\) and exp is étale, then, for all \(x \in \mathrm{L}(\mathrm{G})\) , \(\operatorname{ad} x\) has no eigenvalue belonging to \(2i\pi (\mathbf{Z} - \{0\})\) (§ 6, no. 4, Corollary to Proposition 12). Applying this to \(\lambda x\) , where \(\lambda\) varies through \(\mathbf{C}\) , it follows that all the eigenvalues of ad \(x\) are zero and hence that ad \(x\) is nilpotent. Therefore \(\mathrm{L}(\mathrm{G})\) is nilpotent (Chapter I, § 4, Corollary 1 to Theorem 1) and hence G is nilpotent (Proposition 12).

PROPOSITION 15. Let \(\mathbf{G}\) be a finite-dimensional connected nilpotent Lie group over \(\mathbf{R}\) or \(\mathbf{C}\) and \(\mathbf{A}\) an integral subgroup of \(\mathbf{G}\) . Then \(Z_{\mathrm{G}}(\mathbf{A})\) is the connected Lie subgroup of \(\mathbf{G}\) with Lie algebra \(\mathfrak{z}_{\mathfrak{g}}(\mathrm{L}(\mathbf{A}))\) .

By Proposition 9 of no. 3 it suffices to prove that \(Z_{\mathrm{G}}(\mathbf{A})\) is connected. Let \(g \in Z_{\mathrm{G}}(\mathbf{A})\) . There exists \(x \in \mathrm{L}(\mathbf{G})\) such that \(g = \exp x\) (Proposition 14). Then \(\operatorname{Ad} g|\mathrm{L}(\mathbf{A}) = 1\) (no. 3, Proposition 9), hence \(\operatorname{Ad} g^n |\mathrm{L}(\mathbf{A}) = 1\) for all \(n \in \mathbf{Z}\) and hence \(\exp (\operatorname{ad} nx)|\mathrm{L}(\mathbf{A}) = 1\) for all \(n \in \mathbf{Z}\) . As the mapping

\[
\lambda \mapsto \exp (\operatorname{ad} \lambda x) | \mathrm{L} (\mathrm{A})
\]

of K into \(\mathcal{L}(\mathrm{L}(\mathrm{A}),\mathrm{L}(\mathrm{G}))\) is polynomial, \(\exp(\operatorname{ad}\lambda x)|\mathrm{L}(\mathrm{A})=1\) for all \(\lambda\in K\) , that is \(\exp(\lambda x)\in Z_{\mathrm{G}}(\mathrm{A})\) for all \(\lambda\in K\) .

PROPOSITION 16. Let G be a finite-dimensional nilpotent Lie group over R or C and A an integral subgroup of G distinct from G. Then \(\mathrm{N}_{\mathrm{G}}(\mathrm{A})\) is a connected Lie subgroup of G distinct from A.

\(N_{\mathrm{G}}(\mathsf{A}) \neq \mathsf{A}\) (Algebra, Chapter I, § 6, Corollary 1 to Proposition 8). By Proposition 11 of no. 4, we need only prove that \(N_{\mathrm{G}}(\mathsf{A})\) is connected. Let \(g \in N_{\mathrm{G}}(\mathsf{A})\) . There exists \(x \in \mathrm{L}(\mathrm{G})\) such that \(g = \exp x\) (Proposition 14). Let \(\mathrm{E}\) be the vector subspace of \(\mathcal{L}(\mathrm{L}(\mathrm{G}))\) consisting of the \(u \in \mathcal{L}(\mathrm{L}(\mathrm{G}))\) such that \(u(\mathrm{L}(\mathsf{A})) \subset \mathrm{L}(\mathsf{A})\) . Then \(\operatorname{Ad} g^n \in \mathbb{E}\) and hence \(\exp (\operatorname{ad} nx) \in \mathbb{E}\) for all \(n \in \mathbf{Z}\) . Hence \(\exp (\operatorname{ad} \lambda x) \in \mathbb{E}\) for all \(\lambda \in K\) , that is \(\exp (\lambda x) \in N_{\mathrm{G}}(\mathsf{A})\) for all \(\lambda \in K\) .

PROPOSITION 17. Let \(\mathfrak{g}\) be a finite-dimensional nilpotent Lie algebra over \(\mathbf{K}\) and \((\mathfrak{g}_0, \mathfrak{g}_1, \ldots, \mathfrak{g}_n)\) a decreasing sequence of ideals of \(\mathfrak{g}\) such that \(\mathfrak{g}_0 = \mathfrak{g}, \mathfrak{g}_n = \{0\}\) and \([\mathfrak{g}, \mathfrak{g}_i] \subset \mathfrak{g}_{i+1}\) for \(0 \leqslant i < n\) . Let \(\mathfrak{a}_1, \mathfrak{a}_2, \ldots, \mathfrak{a}_p\) be vector subspaces of \(\mathfrak{g}\) such that each \(\mathfrak{g}_i\) is the direct sum of its intersections with the \(\mathfrak{a}_j\) . Let \(\mathfrak{g}\) be given the Hausdorff law of composition H. Let \(\phi\) be the mapping

\[
\left(x _ {1}, x _ {2}, \dots , x _ {p}\right) \mapsto x _ {1} \mathsf {H} x _ {2} \mathsf {H} \dots \mathsf {H} x _ {p}
\]

\[
o f a _ {1} \times a _ {2} \times \dots \times a _ {p}
\]

\begin{enumerate}
\def\labelenumi{(\roman{enumi})}
\item
  \(\phi\) is a bijection of \(a_{1} \times a_{2} \times \cdots \times a_{p}\) onto \(g\) ;
\item
  \(\phi\) and \(\phi^{-1}\) are polynomial mappings;
\item
  the mapping \((x,y)\mapsto\phi^{-1}(\phi(x)\cdot\phi(y)^{-1})\) of \((\mathfrak{a}_{1}\times\mathfrak{a}_{2}\times\cdots\times\mathfrak{a}_{p})^{2}\) into \(a_{1}\times a_{2}\times\cdots\times a_{p}\) is polynomial.
\end{enumerate}

The proposition is obvious for \(\dim \mathfrak{g} = 0\) . Suppose that \(\dim \mathfrak{g} > 0\) and that the proposition has been established for dimensions \(<\dim \mathfrak{g}\) . It can be assumed that \(\mathfrak{g}_{n-1} \neq \{0\}\) and \(\mathfrak{g}_{n-1}\) is then a non-zero central ideal of \(\mathfrak{g}\) . There exists an index \(j\) such that \(\mathfrak{b} = \mathfrak{g}_{n-1} \cap \mathfrak{a}_j \neq \{0\}\) . Let \(\mathfrak{g}' = \mathfrak{g}/\mathfrak{b}\) , \(\theta\) be the canonical morphism of \(\mathfrak{g}\) onto \(\mathfrak{g}'\) , \(\mathfrak{g}'_i = \theta(\mathfrak{g}_i)\) and \(\mathfrak{a}_i' = \theta(\mathfrak{a}_i)\) . Then \((\mathfrak{g}_0', \mathfrak{g}_1', \ldots, \mathfrak{g}_n')\) is a decreasing sequence of ideals of \(\mathfrak{g}'\) such that \(\mathfrak{g}_0' = \mathfrak{g}'\) , \(\mathfrak{g}_n' = \{0\}\) ,

\[
[ \mathfrak {g} ^ {\prime}, \mathfrak {g} _ {i} ^ {\prime} ] \subset \mathfrak {g} _ {i + 1} ^ {\prime}
\]

for \(0 \leqslant i < n\) and each \(\mathfrak{g}_i'\) is the direct sum of its intersections with the \(\mathfrak{a}_j'\) . Let \(\phi'\) be the mapping

\[
\left(x _ {1} ^ {\prime}, x _ {2} ^ {\prime}, \dots , x _ {p} ^ {\prime}\right) \mapsto x _ {1} ^ {\prime} \mathsf {H} x _ {2} ^ {\prime} \mathsf {H} \dots \mathsf {H} x _ {p} ^ {\prime}
\]

of \(a_{1}^{\prime} \times a_{2}^{\prime} \times \cdots \times a_{p}^{\prime}\) into \(g^{\prime}\) . By the induction hypothesis, \(\phi^{\prime}\) is bijective and \(\phi^{\prime}, \phi^{\prime-1}\) are polynomial mappings.

Let \(x \in g\) . We write

\[
\phi^ {\prime - 1} (\theta (x)) = \left(x _ {1} ^ {\prime} (x), x _ {2} ^ {\prime} (x), \dots , x _ {p} ^ {\prime} (x)\right).\tag{1}
\]

Then

\[
\theta (x) = x _ {1} ^ {\prime} (x) \mathsf {H} x _ {2} ^ {\prime} (x) \mathsf {H} \dots \mathsf {H} x _ {p} ^ {\prime} (x).\tag{2}
\]

Let \(\mathfrak{h}_1\) be a vector subspace supplementary to \(\mathfrak{b}\) in \(\mathfrak{g}\) , the sum of the \(a_k\) with \(k \neq j\) and a supplement of \(\mathfrak{b}\) in \(a_j\) . There exists a bijection \(\eta\) of \(\mathfrak{g}'\) onto \(\mathfrak{h}_1\) such that \(\theta \circ \eta = \mathrm{Id}_{\mathfrak{g}'}\) . For \(x \in \mathfrak{g}\) , we write

\[
\zeta (x) = \eta (x _ {1} ^ {\prime} (x)) \mathsf {H} \eta (x _ {2} ^ {\prime} (x)) \mathsf {H} \dots \mathsf {H} \eta (x _ {p} ^ {\prime} (x)) \in \mathfrak {g}.\tag{3}
\]

\[
y (x) = \zeta (x) ^ {- 1} \text { H } x = (- \zeta (x))\cdot x.\tag{4}
\]

By (2) and (3), \(\theta (\zeta (x)) = \theta (x)\) and hence \(y(x)\in \mathfrak{h}\) . Finally we write

\[
\psi (x) = \left(\eta \left(x _ {1} ^ {\prime} (x)\right), \dots , \eta \left(x _ {j} ^ {\prime} (x)\right) + y (x), \dots , \eta \left(x _ {p} ^ {\prime} (x)\right)\right) \in a _ {1} \times \dots \times a _ {p}.\tag{5}
\]

As \(y(x)\) is central in \(\mathfrak{g}\) ,

\[
\begin{array}{l l} \phi (\psi (x)) = \eta (x _ {1} ^ {\prime} (x)) \texttt {H} \dots \texttt {H} \eta (x _ {j} ^ {\prime} (x)) \texttt {H} \dots \texttt {H} \eta (x _ {p} ^ {\prime} (x)) \texttt {H} y (x) \\ = \zeta (x) \texttt {H} y (x) & \text {by (3)} \\ = x & \text {by (4)}. \end{array}
\]

Hence \(\phi \circ \psi = \mathrm{Id}_{\mathfrak{g}'}\) . Now let \((x_{1}, x_{2}, \ldots, x_{p}) \in \mathfrak{a}_{1} \times \mathfrak{a}_{2} \times \dots \times \mathfrak{a}_{p}\) and write \(x = \phi(x_{1}, x_{2}, \ldots, x_{p}) = x_{1} \mathsf{H} x_{2} \mathsf{H} \dots \mathsf{H} x_{p}\) . Then

\[
\theta (x) = \theta (x _ {1}) \mathsf {H} \theta (x _ {2}) \mathsf {H} \dots \mathsf {H} \theta (x _ {p}),
\]

hence \(x_{i}^{\prime}(x) = \theta (x_{i})\) for \(1\leqslant i\leqslant p\) and therefore

\[
\begin{array}{c} \zeta (x) = x _ {1} \mathsf {H} x _ {2} \mathsf {H} \dots \mathsf {H} (\eta \theta (x _ {j})) \mathsf {H} \dots \mathsf {H} x _ {p} \\ y (x) = x _ {j} - \eta \theta (x _ {j}). \end{array}
\]

Then by (5)

\[
\psi (x) = \left(x _ {1}, \dots , \eta \theta \left(x _ {j}\right) + x _ {j} - \eta \theta \left(x _ {j}\right), \dots , x _ {p}\right) = \left(x _ {1}, x _ {2}, \dots , x _ {p}\right).
\]

Hence \(\psi \circ \phi = \mathrm{Id}_{\mathfrak{a}_1 \times \ldots \times \mathfrak{a}_p}\) . This proves (i). As the Hausdorff law is polynomial, \(\phi\) is polynomial. By the induction hypothesis, \(\phi'^{-1}\) is polynomial; by formula (1), the functions \(x_j'\) are polynomial, hence \(\zeta\) is polynomial (formula (3)), \(y\) is polynomial (formula (4)) and \(\psi\) is polynomial (formula (5)). This proves (ii). Assertion (iii) follows from (i) and (ii) and the fact that the Hausdorff law is polynomial.

Example of a nilpotent Lie group. Let G be the lower strict triangular subgroup of \(\mathbf{GL}(n,\mathbf{K})\) . It is a Lie subgroup of \(\mathbf{GL}(n,\mathbf{K})\) and \(\mathrm{L}(\mathrm{G})\subset\mathrm{gl}(n,\mathbf{K})\) is the Lie algebra of lower triangular matrices with zero diagonal (§3, no. 10, Proposition 36). By Chapter II, §4, no. 6, Remark, G is nilpotent. Suppose henceforth that K=R or C. As G is homeomorphic to \(\mathbf{K}^{n(n-1)/2}\) , G is simply connected. The exponential mapping of L(G) into G is just the mapping

\[
u \mapsto \exp u = \sum_ {k \geqslant 0} \frac {u ^ {k}}{k !} = \sum_ {k = 0} ^ {n - 1} \frac {u ^ {k}}{k !}
\]

(§ 6, no. 4, Example). By Proposition 13, the exponential is an isomorphism of the manifold L(G) onto the manifold G. Proposition 17 of § 6, no. 9 gives the inverse bijection log. We give \(\mathbf{K}^{n}\) a norm. By Spectral Theories, Chapter I, § 4, no. 9, for g ∈ G and \textbar\textbar g − 1\textbar\textbar{} \textless{} 1,

\[
\log g = \sum_ {k \geqslant 1} \frac {(- 1) ^ {k - 1}}{k} (g - 1) ^ {k}
\]

that is

\[
\log g = \sum_ {k = 1} ^ {n - 1} \frac {(- 1) ^ {k - 1}}{k} (g - 1) ^ {k}.\tag{6}
\]

But the two sides of (6) are analytic functions of \(g\) for \(g \in \mathbf{G}\) and are therefore equal for all \(g \in \mathbf{G}\) .

PROPOSITION 18. Let \(k\) be a commutative field. V a vector space of finite dimension \(>0\) over \(k\) and G a subgroup of \(\mathbf{GL}(\mathbf{V})\) whose elements are unipotent.

\begin{enumerate}
\def\labelenumi{(\roman{enumi})}
\item
  There exists a non-zero element \(v\) of \(V\) such that \(gv = v\) for all \(g \in G\) .
\item
  There exists a basis B of V such that, for all \(g \in G\) , the matrix of \(g\) with respect to B is lower triangular and has all its diagonal elements equal to 1.
\item
  The group G is nilpotent.
\end{enumerate}

\begin{enumerate}
\def\labelenumi{(\alph{enumi})}
\tightlist
\item
  Suppose first that \(k\) is algebraically closed and that the identity representation of \(G\) is simple. Let \(a, b\) be in \(G\) . Then
\end{enumerate}

\[
\operatorname{Tr} (a (b - 1)) = \operatorname{Tr} (a b - 1) - \operatorname{Tr} (a - 1) = 0 - 0 = 0
\]

for \(ab - 1\) and \(a - 1\) are nilpotent. As the vector subspace of \(\mathcal{L}(\mathrm{V})\) generated by \(\mathbf{G}\) is \(\mathcal{L}(\mathrm{V})\) (Algebra, Chapter VIII, §4, Corollary 1 to Proposition 2), \(\operatorname{Tr}(u(b - 1)) = 0\) for all \(u \in \mathcal{L}(\mathrm{V})\) and hence \(b = 1\) . Thus \(\mathrm{G} = \{1\}\) .

\begin{enumerate}
\def\labelenumi{(\alph{enumi})}
\setcounter{enumi}{1}
\item
  We pass now to the general case. Let \(\overline{k}\) be an algebraic closure of \(k\) , \(\overline{\mathrm{V}} = \mathrm{V} \otimes_k \overline{k}\) and \(\overline{\mathbf{G}} \subset \mathbf{GL}(\overline{\mathrm{V}})\) the set of \(a \otimes 1\) for \(a \in \mathbf{G}\) . Let \(\mathrm{W}\) (resp. \(\mathrm{W}'\) ) be the set of elements of \(\mathrm{V}\) (resp. \(\overline{\mathrm{V}}\) ) invariant under \(\mathrm{G}\) (resp. \(\overline{\mathbf{G}}\) ). Then \(\mathrm{W}' = \mathrm{W} \otimes_k \overline{k}\) for \(\mathrm{W} = \bigcap_{g \in \mathbb{G}} \mathrm{Ker}(g - 1)\) and \(\mathrm{W}' = \bigcap_{g \in \mathbb{G}} \mathrm{Ker}(g - 1) \otimes 1\) . If \(\mathrm{V}_1\) denotes a minimal element in the set of non-zero vector subspaces of \(\overline{\mathrm{V}}\) which are stable under \(\overline{\mathbf{G}}\) , then \(\mathrm{V}_1 \subset \mathrm{W}'\) by part (a) of the proof; hence \(\mathrm{W} \neq \{0\}\) , which proves (i).
\item
  By induction on dim V, it follows from (i) that there exists an increasing sequence \(\left(\mathrm{V}_{1},\mathrm{V}_{2},\ldots,\mathrm{V}_{n}\right)\) of vector subspaces of V which are stable under G such that \(V_{n}=V\) and the automorphism group of \(V_{i}/V_{i-1}\) canonically derived from G reduces to \(\{1\}\) for all i (we make the convention that \(V_{r}=\{0\}\) for \(r\leqslant0\) ). This implies first (ii) and consequently (iii) (Chapter II, §4, no. 6, Remark).
\end{enumerate}

COROLLARY 1. Let G be a finite-dimensional connected real or complex Lie group. For G to be nilpotent, it is necessary and sufficient that every element of Ad G be nilpotent.

If every element of Ad G is unipotent, Ad G is nilpotent (Proposition 18) and hence G, which is a central extension of Ad G, is nilpotent. If G is nilpotent, L(G) is nilpotent, hence ad x is nilpotent for all \(x \in \mathrm{L}(G)\) and hence \(\mathrm{Ad}(\exp x) = \exp \mathrm{ad} x\) is unipotent; but every element of G is of the form \(\exp x\) for some \(x \in \mathrm{L}(G)\) (Proposition 14).

COROLLARY 2. Every analytic subgroup of \(\mathbf{GL}(n, \mathbf{K})\) consisting of unipotent elements is a simply connected Lie subgroup.

This follows from Propositions 13 (ii) and 18 (ii) and the fact that the lower strict triangular group is simply connected.

